\section{Η Γέννηση του \ENG{ADL} και του \ENG{SCORM} (1997-2004)}
\label{sec:adl_scorm_birth}

Η ανάπτυξη του \textit{\ENG{Advanced} \ENG{Distributed} \ENG{Learning}} (\ENG{ADL}) \ENG{initiative} και του \textit{\ENG{Sharable} \ENG{Content} \ENG{Object} \ENG{Reference} \ENG{Model}} (\ENG{SCORM}) αποτελεί την πιο σημαντική καμπή στην ιστορία των \ENG{E-Learning} \ENG{standards.} Το \ENG{SCORM} έγινε το \ENG{de} \ENG{facto} πρότυπο για την ανάπτυξη, συσκευασία και διανομή εκπαιδευτικού περιεχομένου παγκοσμίως, επηρεάζοντας βαθιά τον τρόπο που σχεδιάζονται και υλοποιούνται τα σύγχρονα συστήματα \ENG{E-Learning}~\cite{adl2011scorm_guide,wiki_adl}.

\subsection{Το \ENG{Advanced} \ENG{Distributed} \ENG{Learning} \ENG{Initiative}}
\label{subsec:adl_initiative}

\subsubsection{Ίδρυση και Νομοθετικό Πλαίσιο}

Το \ENG{ADL} \ENG{initiative} ιδρύθηκε το 1997 από το Υπουργείο Άμυνας των Ηνωμένων Πολιτειών (\ENG{Department} \ENG{of} \ENG{Defense}, \ENG{DoD}) με στόχο την εκσυγχρονισμό της στρατιωτικής εκπαίδευσης μέσω της αξιοποίησης τεχνολογιών κατανεμημένης μάθησης (\textit{\ENG{distributed} \ENG{learning}})~\cite{dod1999strategic,dod2000implementation}. Η δημιουργία του \ENG{ADL} ήταν άμεσο αποτέλεσμα της \textit{\ENG{Quadrennial} \ENG{Defense} \ENG{Review}} του 1996, η οποία αναγνώρισε την ανάγκη για μετασχηματισμό των εκπαιδευτικών συστημάτων των ενόπλων δυνάμεων προκειμένου να ανταποκριθούν στις απαιτήσεις του 21ου αιώνα~\cite{teamorlando2023adl}.


\begin{table}[H]
\centering
\caption{Νομοθετικές πράξεις που θεμελίωσαν το \ENG{ADL}}
\label{tab:adl_legislation}
\small
\setlength{\tabcolsep}{5pt}
\renewcommand{\arraystretch}{1.25}
\begin{chaptertablebox}
\begin{tabular}{|p{5cm}|p{8.5cm}|}
\hline
\textbf{Πράξη} & \textbf{Συμβολή} \\
\hline
\hline
\textbf{\ENG{Public} \ENG{Law} 105-261} (\ENG{Strom} \ENG{Thurmond} \ENG{National} \ENG{Defense} \ENG{Authorization} \ENG{Act}, 1999) & Παρείχε νομική βάση και χρηματοδότηση για την ανάπτυξη κοινών τεχνολογικών προτύπων εκπαίδευσης στο \ENG{DoD}~\cite{congress1999ndaa}. \\
\hline
\textbf{\ENG{Executive} \ENG{Order} 13111} (1999) & Επέκτεινε την εφαρμογή του \ENG{distributed} \ENG{learning} σε ολόκληρη την ομοσπονδιακή κυβέρνηση, καθιστώντας το \ENG{ADL} πανομοσπονδιακό πρόγραμμα~\cite{whitehouse1999eo13111}. \\
\hline
\textbf{\ENG{Department} \ENG{of} \ENG{Defense} \ENG{Strategic} \ENG{Plan} \ENG{for} \ENG{ADL}} (1999) & Όρισε τους στρατηγικούς στόχους: διαλειτουργικότητα, αποθετήρια επαναχρησιμοποιήσιμου περιεχομένου και ενσωμάτωση παιδαγωγικών προσεγγίσεων~\cite{dod1999strategic}. \\
\hline
\end{tabular}
\end{chaptertablebox}
\end{table}


Το \ENG{ADL} λειτουργεί υπό την εποπτεία της \textit{\ENG{Defense} \ENG{Human} \ENG{Resources} \ENG{Activity}} (\ENG{DHRA}) και έχει ως αποστολή την προώθηση της καινοτομίας στην εκπαίδευση και την κατάρτιση μέσω της έρευνας, της ανάπτυξης προτύπων και της συνεργασίας με ακαδημαϊκούς και βιομηχανικούς φορείς.

\subsubsection{Αποστολή και Πεδία Εστίασης}

Το \ENG{ADL} επικεντρώνεται σε πέντε κύριους τομείς: \ENG{E-Learning} \ENG{Standards}, με έμφαση σε τεχνικές προδιαγραφές διαλειτουργικότητας (\ENG{SCORM}, \ENG{xAPI}); \ENG{Mobile} \ENG{Learning}, με λύσεις για μάθηση μέσω κινητών συσκευών; \ENG{Simulations} \ENG{and} \ENG{Virtual} \ENG{Worlds}, με ενσωμάτωση προσομοιώσεων, εικονικών περιβαλλόντων και \textit{\ENG{serious} \ENG{games}} στην εκπαίδευση; \ENG{Learning} \ENG{Analytics}, με αξιοποίηση δεδομένων για αξιολόγηση και βελτίωση της μαθησιακής εμπειρίας· και \ENG{Instructional} \ENG{Design} \ENG{and} \ENG{Learning} \ENG{Theory}, με έρευνα σε παιδαγωγικά μοντέλα και θεωρίες μάθησης.

Το \ENG{ADL} συνεργάζεται με πανεπιστήμια, ερευνητικά κέντρα και εταιρείες τεχνολογίας για την προώθηση της έρευνας και την πιλοτική εφαρμογή νέων τεχνολογιών. Μέσω αυτών των συνεργασιών, το \ENG{ADL} έχει καταστεί παγκόσμιος ηγέτης στην προτυποποίηση του \ENG{E-Learning.}

\subsection{\ENG{SCORM} 1.0 και 1.1: Οι Πρώτες Προσπάθειες (2000)}
\label{subsec:scorm_10_11}

Η πρώτη έκδοση του \ENG{SCORM} (1.0) κυκλοφόρησε τον Ιανουάριο του 2000 και ακολουθήθηκε από τη βελτιωμένη έκδοση 1.1 τον Ιανουάριο του 2001. Αυτές οι αρχικές εκδόσεις αποτελούσαν πρωτόγονες προσπάθειες συγκέντρωσης και ενοποίησης υπαρχόντων προτύπων και προδιαγραφών από διάφορους οργανισμούς, συμπεριλαμβανομένων του \ENG{AICC}, του \ENG{IMS} \ENG{Global} \ENG{Learning} \ENG{Consortium} και του \ENG{IEEE}~\cite{scormcom2024versions}.

Το \ENG{SCORM} 1.0/1.1 εισήγαγε τρεις βασικές έννοιες: το \ENG{Content} \ENG{Aggregation} \ENG{Model} (\ENG{CAM}) ως μοντέλο συσκευασίας και οργάνωσης περιεχομένου (συνδυασμός \ENG{assets} και \ENG{SCOs} σε πακέτα διανομής), το \ENG{Run-Time} \ENG{Environment} (\ENG{RTE}) ως προδιαγραφές επικοινωνίας περιεχομένου-\ENG{LMS} μέσω \ENG{JavaScript} \ENG{API}, και τα \ENG{Metadata} μέσω μιας απλοποιημένης εκδοχής του \ENG{IEEE} \ENG{LOM}.

Παρά τη σημασία τους ως πρώτες προσπάθειες, οι εκδόσεις 1.0 και 1.1 παρουσίαζαν ασαφείς ή ελλιπείς προδιαγραφές σε κρίσιμα σημεία, περιορισμένες δυνατότητες αλληλουχίας (\textit{\ENG{sequencing}}) και ασυνέπειες στην υλοποίηση από διαφορετικούς προμηθευτές \ENG{LMS.}

Αυτές οι αδυναμίες οδήγησαν στην ανάπτυξη μιας πιο ώριμης και σταθερής έκδοσης: το \ENG{SCORM} 1.2.

\subsection{\ENG{SCORM} 1.2: Η Σταθεροποίηση του Προτύπου (2001)}
\label{subsec:scorm_12}

Το \ENG{SCORM} 1.2, που κυκλοφόρησε τον Οκτώβριο του 2001, αποτέλεσε την πρώτη ευρέως αποδεκτή και σταθερή έκδοση του προτύπου~\cite{scormcom2024scorm12overview,adl2002cookbook}. Η έκδοση αυτή έλυσε πολλά από τα προβλήματα των προηγούμενων εκδόσεων και κατέστη το πιο δημοφιλές \ENG{SCORM} πρότυπο για πολλά χρόνια. Ακόμη και σήμερα, πολλά \ENG{LMS} και εκπαιδευτικά εργαλεία συνεχίζουν να υποστηρίζουν \ENG{SCORM} 1.2 περιεχόμενο.

\subsubsection{Βασικά Χαρακτηριστικά του \ENG{SCORM} 1.2}

\paragraph{\ENG{Content} \ENG{Aggregation} \ENG{Model} (\ENG{CAM})}

Το \ENG{SCORM} 1.2 \ENG{CAM} όρισε μια ιεραρχική δομή για την οργάνωση εκπαιδευτικού περιεχομένου.

\begin{table}[H]
\centering
\caption{Βασικά στοιχεία του \ENG{SCORM} 1.2 \ENG{CAM}}
\label{tab:scorm12_cam}
\small
\setlength{\tabcolsep}{5pt}
\renewcommand{\arraystretch}{1.25}
\begin{chaptertablebox}
\begin{tabular}{|p{4.2cm}|p{9.3cm}|}
\hline
\textbf{Στοιχείο} & \textbf{Περιγραφή} \\
\hline
\hline
\textbf{\ENG{Assets}} & Βασικά δομικά στοιχεία (αρχεία \ENG{HTML}, εικόνες, βίντεο, \ENG{JavaScript}). Δεν επικοινωνούν με το \ENG{LMS.} \\
\hline
\textbf{\ENG{Sharable} \ENG{Content} \ENG{Objects} (\ENG{SCOs})} & Μονάδες περιεχομένου που επικοινωνούν με το \ENG{LMS} μέσω \ENG{SCORM} \ENG{API} και αποτελούν την ελάχιστη μονάδα \textit{\ENG{tracking}}. \\
\hline
\textbf{\ENG{Content} \ENG{Aggregations}} & Συλλογές από \ENG{SCOs} και \ENG{assets} οργανωμένες σε λογική δομή (π.χ. μάθημα, ενότητα, κεφάλαιο). \\
\hline
\textbf{\ENG{imsmanifest.xml}} & \ENG{XML} αρχείο που περιγράφει τη δομή, τα \ENG{metadata}, τις εξαρτήσεις και τη σχέση μεταξύ \ENG{SCOs} και \ENG{assets}. Λειτουργεί ως «χάρτης» για το \ENG{LMS.} \\
\hline
\end{tabular}
\end{chaptertablebox}
\end{table}


Η δομή του \texttt{\ENG{imsmanifest.xml}} περιλάμβανε τις εξής βασικές ενότητες:

{\selectlanguage{english}\fontencoding{T1}\selectfont

\begin{verbatim}
<manifest>
    <metadata>...</metadata>
    <organizations>
        <organization>
            <item identifier="item_1" identifierref="resource_1">
                <title>Εισαγωγή</title>
            </item>
            ...
        </organization>
    </organizations>
    <resources>
        <resource identifier="resource_1" type="webcontent" 
                  href="intro.html">
            <file href="intro.html"/>
            <file href="images/logo.png"/>
        </resource>
        ...
    </resources>
</manifest>
\end{verbatim}
}
\paragraph{\ENG{Run-Time} \ENG{Environment} (\ENG{RTE})}

Το \ENG{SCORM} 1.2 \ENG{RTE} όριζε τις βασικές κλήσεις επικοινωνίας που παρουσιάζονται στον παρακάτω πίνακα.

\begin{table}[H]
\centering
\caption{Βασικές λειτουργίες του \ENG{SCORM} 1.2 \ENG{RTE}}
\label{tab:scorm12_rte_functions}
\small
\setlength{\tabcolsep}{5pt}
\renewcommand{\arraystretch}{1.25}
\begin{chaptertablebox}
\begin{tabular}{|p{5.2cm}|p{8.3cm}|}
\hline
\textbf{Λειτουργία} & \textbf{Ρόλος} \\
\hline
\hline
\texttt{\ENG{LMSInitialize}()} & Εκκινεί το \ENG{SCO} και εγκαθιδρύει επικοινωνία με το \ENG{LMS.} \\
\hline
\texttt{\ENG{LMSFinish}()} & Ολοκληρώνει το \ENG{SCO} και κλείνει την επικοινωνία. \\
\hline
\texttt{\ENG{LMSGetValue}(\ENG{element})} & Ανακτά την τιμή ενός στοιχείου από το \ENG{LMS.} \\
\hline
\texttt{\ENG{LMSSetValue}(\ENG{element}, \ENG{value})} & Ορίζει τιμή σε στοιχείο του \ENG{LMS.} \\
\hline
\texttt{\ENG{LMSCommit}()} & Αποθηκεύει άμεσα τα δεδομένα στο \ENG{LMS} (προαιρετικό). \\
\hline
\texttt{\ENG{LMSGetLastError}()} & Επιστρέφει τον τελευταίο κωδικό σφάλματος. \\
\hline
\texttt{\ENG{LMSGetErrorString}(\ENG{errorCode})} & Παρέχει περιγραφή για τον κωδικό σφάλματος. \\
\hline
\texttt{\ENG{LMSGetDiagnostic}(\ENG{errorCode})} & Παρέχει διαγνωστικές πληροφορίες για το σφάλμα. \\
\hline
\end{tabular}
\end{chaptertablebox}
\end{table}


\paragraph{\ENG{Data} \ENG{Model}}

Το \ENG{SCORM} 1.2 \ENG{data} \ENG{model}, βασισμένο στο \ENG{AICC} \ENG{CMI} \ENG{data} \ENG{model}, περιλάμβανε βασικά στοιχεία που συνοψίζονται στον παρακάτω πίνακα.

\begin{table}[H]
\centering
\caption{Βασικά στοιχεία του \ENG{SCORM} 1.2 \ENG{data} \ENG{model}}
\label{tab:scorm12_data_model}
\small
\setlength{\tabcolsep}{5pt}
\renewcommand{\arraystretch}{1.25}
\begin{chaptertablebox}
\begin{tabular}{|p{5cm}|p{8.5cm}|}
\hline
\textbf{Στοιχείο} & \textbf{Περιγραφή} \\
\hline
\hline
\texttt{\ENG{cmi.core.student}\_\ENG{id}} & Το μοναδικό αναγνωριστικό του εκπαιδευόμενου. \\
\hline
\texttt{\ENG{cmi.core.student}\_\ENG{name}} & Το ονοματεπώνυμο του εκπαιδευόμενου. \\
\hline
\texttt{\ENG{cmi.core.lesson}\_\ENG{location}} & Η τελευταία θέση στο \ENG{SCO} (για \ENG{bookmarking}). \\
\hline
\texttt{\ENG{cmi.core.lesson}\_\ENG{status}} & Η κατάσταση του μαθήματος (π.χ. \texttt{\ENG{incomplete}}, \texttt{\ENG{completed}}, \texttt{\ENG{passed}}, \texttt{\ENG{failed}}). \\
\hline
\texttt{\ENG{cmi.core.score.raw}} & Η ακατέργαστη βαθμολογία (αριθμός). \\
\hline
\texttt{\ENG{cmi.core.score.min}/\ENG{max}} & Οι ελάχιστες και μέγιστες δυνατές βαθμολογίες. \\
\hline
\texttt{\ENG{cmi.core.session}\_\ENG{time}} & Ο χρόνος που πέρασε ο εκπαιδευόμενος στο τρέχον \ENG{session.} \\
\hline
\texttt{\ENG{cmi.core.total}\_\ENG{time}} & Ο συνολικός χρόνος σε όλα τα \ENG{sessions.} \\
\hline
\texttt{\ENG{cmi.suspend}\_\ENG{data}} & \ENG{String} για αποθήκευση προσαρμοσμένων δεδομένων (έως 4096 χαρακτήρες). \\
\hline
\texttt{\ENG{cmi.launch}\_\ENG{data}} & Δεδομένα που το \ENG{LMS} περνά στο \ENG{SCO} κατά την εκκίνηση. \\
\hline
\texttt{\ENG{cmi.interactions}.*} & Καταγραφή αλληλεπιδράσεων (ερωτήσεις, απαντήσεις). \\
\hline
\texttt{\ENG{cmi.objectives}.*} & Παρακολούθηση μαθησιακών στόχων. \\
\hline
\end{tabular}
\end{chaptertablebox}
\end{table}


\subsubsection{Ευρεία Αποδοχή και Επιτυχία}

Το \ENG{SCORM} 1.2 γρήγορα έγινε το βιομηχανικό πρότυπο. Η επιτυχία του αποδίδεται στην απλότητα υλοποίησης, στη σταθερότητα και την καθαρή τεκμηρίωση, στην ισχυρή υποστήριξη από προμηθευτές \ENG{LMS} και \ENG{authoring} \ENG{tools}, καθώς και στη διαμόρφωση ενεργής κοινότητας προγραμματιστών, εκπαιδευτών και παρόχων περιεχομένου.

\subsection{\ENG{SCORM} 2004: Η Εξέλιξη προς Προηγμένη Αλληλουχία (2004-2009)}
\label{subsec:scorm_2004}

Παρά την επιτυχία του \ENG{SCORM} 1.2, το \ENG{ADL} αναγνώρισε την ανάγκη για πιο προηγμένες δυνατότητες, ιδίως στον τομέα της αλληλουχίας και πλοήγησης. Το αποτέλεσμα ήταν το \ENG{SCORM} 2004, το οποίο κυκλοφόρησε σε τέσσερις εκδόσεις (\ENG{1st}, \ENG{2nd}, \ENG{3rd}, \ENG{4th} \ENG{Edition}) μεταξύ 2004 και 2009~\cite{adl2009scorm2004,adl2009impact}.

\subsubsection{Νέα Χαρακτηριστικά του \ENG{SCORM} 2004}

\paragraph{\ENG{Sequencing} \ENG{and} \ENG{Navigation} (\ENG{SN})}

Η σημαντικότερη προσθήκη στο \ENG{SCORM} 2004 ήταν το \ENG{Sequencing} \ENG{and} \ENG{Navigation} (\ENG{SN}) \ENG{book}, το οποίο όριζε προηγμένες δυνατότητες για τον έλεγχο της σειράς εμφάνισης των εκπαιδευτικών μονάδων.
\begin{table}[H]
\centering
\caption{Κύρια χαρακτηριστικά \ENG{Sequencing} \& \ENG{Navigation} στο \ENG{SCORM} 2004}
\label{tab:scorm2004_sn_features}
\small
\setlength{\tabcolsep}{5pt}
\renewcommand{\arraystretch}{1.25}
\begin{chaptertablebox}
\begin{tabular}{|p{4.3cm}|p{9.2cm}|}
\hline
\textbf{Χαρακτηριστικό} & \textbf{Περιγραφή} \\
\hline
\hline
\textbf{\ENG{Activity} \ENG{Trees}} & Ιεραρχική οργάνωση εκπαιδευτικών δραστηριοτήτων σε δενδρική δομή. \\
\hline
\textbf{\ENG{Sequencing} \ENG{Rules}} & Κανόνες προϋποθέσεων και μετα-συνθηκών για πρόσβαση σε δραστηριότητες. \\
\hline
\textbf{\ENG{Rollup} \ENG{Rules}} & Κανόνες που καθορίζουν πώς η \ENG{completion/success} των θυγατρικών επηρεάζει την κατάσταση της γονικής δραστηριότητας. \\
\hline
\textbf{\ENG{Navigation} \ENG{Requests}} & Δυνατότητα αιτημάτων πλοήγησης (π.χ. «επόμενο», «προηγούμενο», «πήγαινε σε συγκεκριμένη δραστηριότητα»). \\
\hline
\textbf{\ENG{Control} \ENG{Modes}} & Ρυθμίσεις που καθορίζουν τον τρόπο πλοήγησης (π.χ. \texttt{\ENG{choice}}, \texttt{\ENG{flow}}). \\
\hline
\end{tabular}
\end{chaptertablebox}
\end{table}


\paragraph{Βελτιωμένο \ENG{Data} \ENG{Model}}

Το \ENG{SCORM} 2004 εισήγαγε ένα επεκταμένο \ENG{data} \ENG{model}: αλλαγή προθέματος από \texttt{\ENG{cmi.core.}*} σε \texttt{\ENG{cmi.}*} (π.χ. \texttt{\ENG{cmi.learner}\_\ENG{id}} αντί για \texttt{\ENG{cmi.core.student}\_\ENG{id}}), νέα στοιχεία όπως \texttt{\ENG{cmi.completion}\_\ENG{status}}, βελτιωμένη υποστήριξη στόχων μέσω \texttt{\ENG{cmi.objectives}} με δυνατότητα \textit{\ENG{global} \ENG{objectives}}, και πιο δομημένη καταγραφή αλληλεπιδράσεων μέσω \texttt{\ENG{cmi.interactions}}.

\paragraph{Βελτιωμένο \ENG{API}}

Το \ENG{SCORM} 2004 άλλαξε τα ονόματα των \ENG{API} μεθόδων, όπως συνοψίζεται στον παρακάτω πίνακα.

\begin{table}[H]
\centering
\caption{Μετονομασία \ENG{API} μεθόδων από \ENG{SCORM} 1.2 σε \ENG{SCORM} 2004}
\label{tab:scorm2004_api_changes}
\small
\setlength{\tabcolsep}{5pt}
\renewcommand{\arraystretch}{1.25}
\begin{chaptertablebox}
\begin{tabular}{|p{6.2cm}|p{6.2cm}|}
\hline
\textbf{\ENG{SCORM} 1.2} & \textbf{\ENG{SCORM} 2004} \\
\hline
\hline
\texttt{\ENG{LMSInitialize}()} & \texttt{\ENG{Initialize}()} \\
\hline
\texttt{\ENG{LMSFinish}()} & \texttt{\ENG{Terminate}()} \\
\hline
\texttt{\ENG{LMSGetValue}()} & \texttt{\ENG{GetValue}()} \\
\hline
\texttt{\ENG{LMSSetValue}()} & \texttt{\ENG{SetValue}()} \\
\hline
\texttt{\ENG{LMSCommit}()} & \texttt{\ENG{Commit}()} \\
\hline
\texttt{\ENG{LMSGetLastError}()} & \texttt{\ENG{GetLastError}()} \\
\hline
\end{tabular}
\end{chaptertablebox}
\end{table}


\subsubsection{Οι Τέσσερις Εκδόσεις (\ENG{Editions})}

Το \ENG{SCORM} 2004 κυκλοφόρησε σε τέσσερις εκδόσεις: η \ENG{1st} \ENG{Edition} (2004) αποτέλεσε την αρχική έκδοση με πλήρη υποστήριξη \ENG{SN}, η \ENG{2nd} \ENG{Edition} (2004) εισήγαγε διορθώσεις \ENG{bugs} και βελτιώσεις τεκμηρίωσης, η \ENG{3rd} \ENG{Edition} (2006) βελτίωσε τη σαφήνεια των προδιαγραφών και τη συμβατότητα, ενώ η \ENG{4th} \ENG{Edition} (2009) παραμένει η τελική και πιο σταθερή έκδοση του \ENG{SCORM} 2004~\cite{adl2009scorm2004}.

\subsubsection{Προκλήσεις και Περιορισμένη Υιοθέτηση}

Παρά τις προηγμένες δυνατότητες του, το \ENG{SCORM} 2004 δεν έγινε ποτέ τόσο δημοφιλές όσο το \ENG{SCORM} 1.2. Οι λόγοι σχετίζονται με την πολυπλοκότητα του \ENG{Sequencing} \ENG{and} \ENG{Navigation} και τη μερική υλοποίησή του από προμηθευτές \ENG{LMS} και \ENG{authoring} \ENG{tools}, την έλλειψη \ENG{backward} \ENG{compatibility} με το \ENG{SCORM} 1.2 που απαιτούσε μετατροπές περιεχομένου, την περιορισμένη ζήτηση από οργανισμούς που θεωρούσαν επαρκές το \ENG{SCORM} 1.2, καθώς και το μικρότερο οικοσύστημα εργαλείων.

\subsection{Επίδραση του \ENG{SCORM} στην Αγορά \ENG{E-Learning}}
\label{subsec:scorm_market_impact}

Η επιρροή του \ENG{SCORM} στη βιομηχανία \ENG{E-Learning} είναι αδιαμφισβήτητη. Δημιούργησε ένα ευρύ οικοσύστημα προμηθευτών, εργαλείων και υπηρεσιών (π.χ. \ENG{Rustici} \ENG{Software} με το \ENG{SCORM} \ENG{Cloud}, \ENG{Articulate}, \ENG{Adobe}, \ENG{iSpring}), μείωσε το κόστος ανάπτυξης και συντήρησης μέσω διαλειτουργικότητας, επέτρεψε την άμεση ανταλλαγή και αξιοποίηση έτοιμου περιεχομένου μεταξύ οργανισμών και καθιέρωσε \ENG{best} \ENG{practices} για οργάνωση, παρακολούθηση και αξιολόγηση εκπαιδευτικού περιεχομένου.

Ωστόσο, το \ENG{SCORM} είχε περιορισμούς που οδήγησαν στην ανάπτυξη νεότερων προτύπων όπως το \ENG{xAPI}: ισχυρή εξάρτηση από \ENG{web} \ENG{browsers} που δυσκολεύει \ENG{mobile} \ENG{apps} και \ENG{simulations}, περιορισμένο \ENG{data} \ENG{model} που δεν καταγράφει πλούσιες εμπειρίες, απουσία \ENG{offline} \ENG{tracking} χωρίς σύνδεση με \ENG{LMS}, και μονολιθική δομή γύρω από την έννοια του «μαθήματος», με αδύναμη υποστήριξη \ENG{social} και \ENG{informal} \ENG{learning}.

Παρά αυτούς τους περιορισμούς, το \ENG{SCORM} παραμένει ευρέως χρησιμοποιούμενο και θα συνεχίσει να παίζει σημαντικό ρόλο στο \ENG{E-Learning} για αρκετά χρόνια ακόμη.

% Βιβλιογραφικές αναφορές για ενότητα 2.5
% Προσθήκη στο αρχείο: \ENG{bibliography/chapter2}_\ENG{historical.bib}

% \ENG{Figure}: \ENG{SCORM} \ENG{Architecture}
\begin{figure}[H]
  \centering
  % FIGURE 2.5: SCORM Architecture (3-Tier Model)
% Content Aggregation Model (CAM) + Run-Time Environment (RTE)
\begin{chapterfigurebox}
\begin{tikzpicture}[
  scale=0.9,
  every node/.style={font=\small},
  layer/.style={rectangle, draw, thick, minimum width=12cm, minimum height=1.8cm, align=center, fill=white},
  component/.style={rectangle, rounded corners, draw, fill=white, minimum width=2.5cm, minimum height=1cm, align=center, font=\footnotesize},
  arrow/.style={-latex, thick}
]

% Layer 1: Content Aggregation Model (CAM) - Top
\node[layer, fill=scormgreen!10] at (0, 6) (cam) {
  \textbf{\large \ENG{Content} \ENG{Aggregation} \ENG{Model} (\ENG{CAM})}\\[0.2cm]
  \small Δομή και \ENG{Packaging} Περιεχομένου
};

% CAM Components
\node[component, fill=scormgreen!30] at (-4, 4) (assets) {
  \textbf{\ENG{Assets}}\\
  \tiny \ENG{Media} \ENG{files}\\
  \tiny (\ENG{HTML}, \ENG{images}, \ENG{PDF})
};
\node[component, fill=scormgreen!40] at (-1, 4) (sco) {
  \textbf{\ENG{SCOs}}\\
  \tiny \ENG{Shareable} \ENG{Content}\\
  \tiny \ENG{Objects}
};
\node[component, fill=scormgreen!50] at (2, 4) (activities) {
  \textbf{\ENG{Activities}}\\
  \tiny \ENG{Learning}\\
  \tiny \ENG{Activities}
};
\node[component, fill=scormgreen!60] at (5, 4) (org) {
  \textbf{\ENG{Organizations}}\\
  \tiny \ENG{Content}\\
  \tiny \ENG{Structure}
};

% Arrows between CAM components
\draw[arrow, scormgreen] (assets) -- (sco);
\draw[arrow, scormgreen] (sco) -- (activities);
\draw[arrow, scormgreen] (activities) -- (org);

% Layer 2: Run-Time Environment (RTE) - Middle
\node[layer, fill=blue!10] at (0, 1.5) (rte) {
  \textbf{\large \ENG{Run-Time} \ENG{Environment} (\ENG{RTE})}\\[0.2cm]
  \small Επικοινωνία \ENG{LMS} \ensuremath{\leftrightarrow} \ENG{Content}
};

% RTE Components
\node[component, fill=blue!30] at (-4, -0.3) (launch) {
  \textbf{\ENG{Launch}}\\
  \tiny Εκκίνηση \ENG{SCO}
};
\node[component, fill=blue!40] at (-1, -0.3) (api) {
  \textbf{\ENG{API} \ENG{Adapter}}\\
  \tiny \texttt{\ENG{Initialize}()}\\
  \tiny \texttt{\ENG{Terminate}()}
};
\node[component, fill=blue!50] at (2, -0.3) (datamodel) {
  \textbf{\ENG{Data} \ENG{Model}}\\
  \tiny \ENG{CMI} \ENG{Elements}
};
\node[component, fill=blue!60] at (5, -0.3) (tracking) {
  \textbf{\ENG{Tracking}}\\
  \tiny \ENG{GetValue}()\\
  \tiny \ENG{SetValue}()
};

% Arrows between RTE components
\draw[arrow, blue!80] (launch) -- (api);
\draw[arrow, blue!80] (api) -- (datamodel);
\draw[arrow, blue!80] (datamodel) -- (tracking);

% Layer 3: Sequencing & Navigation (S&N) - Bottom
\node[layer, fill=accentpurple!10] at (0, -3) (sn) {
  \textbf{\large \ENG{Sequencing} \& \ENG{Navigation} (\ENG{S}\&\ENG{N})}\\[0.2cm]
  \small Έλεγχος Ροής και Προϋποθέσεων
};

% S&N Components
\node[component, fill=accentpurple!30] at (-3.5, -4.8) (seq) {
  \textbf{\ENG{Sequencing}}\\
  \tiny \ENG{Rules}, \ENG{Pre/Post-}\\
  \tiny \ENG{conditions}
};
\node[component, fill=accentpurple!40] at (0, -4.8) (nav) {
  \textbf{\ENG{Navigation}}\\
  \tiny \ENG{Continue}, \ENG{Previous},\\
  \tiny \ENG{Choice}
};
\node[component, fill=accentpurple!50] at (3.5, -4.8) (rollup) {
  \textbf{\ENG{Rollup}}\\
  \tiny \ENG{Score} \ENG{aggregation}\\
  \tiny \ENG{Completion} \ENG{status}
};

% Arrows between S&N components
\draw[arrow, accentpurple] (seq) -- (nav);
\draw[arrow, accentpurple] (nav) -- (rollup);

% Vertical connections between layers
\draw[arrow, thick, gray, dashed] (cam.south) -- (rte.north) node[midway, right, font=\tiny] {\ENG{Manifest} (\ENG{imsmanifest.xml})};
\draw[arrow, thick, gray, dashed] (rte.south) -- (sn.north) node[midway, right, font=\tiny] {\ENG{Sequencing} \ENG{rules}};

% Side annotations
\node[right, font=\scriptsize, align=left, text width=3cm] at (7, 4) {
  \textbf{\ENG{SCORM} 1.2:}\\
  \tiny \ENG{Assets} + \ENG{SCOs}\\
  \tiny Βασικό \ENG{RTE}\\
  \tiny Όχι \ENG{S}\&\ENG{N}
};

\node[right, font=\scriptsize, align=left, text width=3cm] at (7, 1.5) {
  \textbf{\ENG{SCORM} 2004:}\\
  \tiny + \ENG{Sequencing}\\
  \tiny + \ENG{Rollup} \ENG{Rules}\\
  \tiny + \ENG{Advanced} \ENG{Nav}
};

% Legend
\node[below, font=\tiny, align=center] at (0, -6.5) {
  \textcolor{scormgreen}{\textbullet} \ENG{CAM}: \ENG{Packaging} \& \ENG{Structure} \quad
  \textcolor{blue}{\textbullet} \ENG{RTE}: \ENG{API} \ENG{Communication} \quad
  \textcolor{accentpurple}{\textbullet} \ENG{S}\&\ENG{N}: \ENG{Flow} \ENG{Control} (\ENG{SCORM} 2004 \ENG{only})
};

\end{tikzpicture}
\end{chapterfigurebox}

  \caption{Αρχιτεκτονική \ENG{SCORM} - \ENG{CAM}, \ENG{RTE}, και \ENG{Sequencing} \& \ENG{Navigation}}
  \label{fig:scorm_architecture}
\end{figure}

Απεικονίζει την τριεπίπεδη αρχιτεκτονική του \ENG{SCORM}, που αποτελείται από το \ENG{Content} \ENG{Aggregation} \ENG{Model} (\ENG{CAM}), το \ENG{Run-Time} \ENG{Environment} (\ENG{RTE}), και το \ENG{Sequencing} \& \ENG{Navigation} (\ENG{S\&N}).

% \ENG{Table}: \ENG{SCORM} \ENG{Versions} \ENG{Evolution}
\setlength{\LTleft}{0pt plus 1fil}
\setlength{\LTright}{0pt plus 1fil}
\small
\setlength{\tabcolsep}{5pt}
\renewcommand{\arraystretch}{1.25}
\begin{longtable}{@{}p{2.2cm}p{1.4cm}p{2.1cm}p{4.0cm}p{4.0cm}@{}}
\caption{Εξέλιξη Εκδόσεων \ENG{SCORM} (2000-2009)}
\label{tab:scorm_versions}\\
\toprule
\textbf{\ENG{Version}} & \textbf{\ENG{Year}} & \textbf{\ENG{Status}} & \textbf{\ENG{Key} \ENG{Features}} & \textbf{\ENG{Limitations}} \\
\midrule
\endfirsthead

\caption[]{Εξέλιξη Εκδόσεων \ENG{SCORM} (2000-2009) (συνέχεια)}\\
\toprule
\textbf{\ENG{Version}} & \textbf{\ENG{Year}} & \textbf{\ENG{Status}} & \textbf{\ENG{Key} \ENG{Features}} & \textbf{\ENG{Limitations}} \\
\midrule
\endhead

\midrule
\multicolumn{5}{r}{\footnotesize Συνεχίζεται στην επόμενη σελίδα}\\
\endfoot

\bottomrule
\endlastfoot

\textbf{\ENG{SCORM} 1.0} 
& \ENG{Jan} 2000 
& \ENG{Obsolete} 
& \textbullet~\ENG{First} \ENG{release}\newline
  \textbullet~\ENG{Basic} \ENG{CAM} (\ENG{Content} \ENG{Aggregation})\newline
  \textbullet~\ENG{AICC-inspired} \ENG{CMI} \ENG{data} \ENG{model}
& \textbullet~\ENG{Immature} \ENG{specification}\newline
  \textbullet~\ENG{Limited} \ENG{adoption}\newline
  \textbullet~\ENG{No} \ENG{sequencing} \\

\midrule

\textbf{\ENG{SCORM} 1.1} 
& \ENG{Jan} 2001 
& \ENG{Obsolete} 
& \textbullet~\ENG{Refined} \ENG{CAM}\newline
  \textbullet~\ENG{Improved} \ENG{metadata} (\ENG{IEEE} \ENG{LOM})\newline
  \textbullet~\ENG{Better} \ENG{API} \ENG{specification}
& \textbullet~\ENG{Still} \ENG{no} \ENG{sequencing}\newline
  \textbullet~\ENG{Compatibility} \ENG{issues} \ENG{with} 1.0 \\

\midrule

\textbf{\ENG{SCORM} 1.2} 
& \ENG{Oct} 2001 
& \ENG{Legacy} (\ENG{widely} \ENG{used})
& \textbullet~\textbf{\ENG{Most} \ENG{adopted} \ENG{version}}\newline
  \textbullet~\ENG{Stable} \ENG{RTE} (\ENG{Run-Time} \ENG{Environment})\newline
  \textbullet~\ENG{Simple} \ENG{API} (\texttt{\ENG{LMSGetValue}})\newline
  \textbullet~\ENG{Manifest-based} \ENG{packaging}\newline
  \textbullet~\ENG{Backward} \ENG{compatible} \ENG{with} \ENG{AICC}
& \textbullet~\ENG{No} \ENG{advanced} \ENG{sequencing}\newline
  \textbullet~\ENG{Limited} \ENG{navigation} \ENG{control}\newline
  \textbullet~\ENG{No} \ENG{rollup} \ENG{rules}\newline
  \textbullet~128-\ENG{char} \ENG{data} \ENG{limits} \\

\midrule

\textbf{\ENG{SCORM} 2004} 
& \ENG{Jan} 2004 
& \ENG{Current} 
& \textbullet~\textbf{\ENG{Simple} \ENG{Sequencing}} (\ENG{IMS} \ENG{SS})\newline
  \textbullet~\ENG{Navigation} \ENG{controls}\newline
  \textbullet~\ENG{Pre/post-conditions}\newline
  \textbullet~\ENG{Rollup} \ENG{rules} (\ENG{score} \ENG{aggregation})\newline
  \textbullet~\ENG{Run-Time} \ENG{Communication} (\ENG{RTC})
& \textbullet~\ENG{High} \ENG{complexity}\newline
  \textbullet~\ENG{Steep} \ENG{learning} \ENG{curve}\newline
  \textbullet~\ENG{Authoring} \ENG{tool} \ENG{challenges}\newline
  \textbullet~\ENG{Not} \ENG{backward} \ENG{compatible} \ENG{with} 1.2 \\

\midrule

\textbf{\ENG{SCORM} 2004 2\ENG{nd} \ENG{Edition}} 
& \ENG{Jul} 2004 
& \ENG{Superseded} 
& \textbullet~\ENG{Bug} \ENG{fixes} \ENG{from} 1\ENG{st} \ENG{edition}\newline
  \textbullet~\ENG{Clarified} \ENG{sequencing} \ENG{behaviors}\newline
  \textbullet~\ENG{Improved} \ENG{conformance} \ENG{tests}
& \textbullet~\ENG{Still} \ENG{complex}\newline
  \textbullet~\ENG{Limited} \ENG{real-world} \ENG{adoption} \\

\midrule

\textbf{\ENG{SCORM} 2004 3\ENG{rd} \ENG{Edition}} 
& \ENG{Oct} 2006 
& \ENG{Superseded} 
& \textbullet~\ENG{Refined} \ENG{sequencing} \ENG{rules}\newline
  \textbullet~\ENG{Better} \ENG{error} \ENG{handling}\newline
  \textbullet~\ENG{Extended} \ENG{API} \ENG{support}\newline
  \textbullet~\ENG{Navigation} \ENG{request} \ENG{improvements}
& \textbullet~\ENG{Complexity} \ENG{remains}\newline
  \textbullet~\ENG{Authoring} \ENG{tools} \ENG{lag} \ENG{behind} \ENG{spec} \\

\midrule

\textbf{\ENG{SCORM} 2004 4\ENG{th} \ENG{Edition}} 
& \ENG{Mar} 2009 
& \textbf{\ENG{Final}} (\ENG{current}) 
& \textbullet~\textbf{\ENG{Current} \ENG{recommended} \ENG{version}}\newline
  \textbullet~\ENG{Mature} \ENG{sequencing} \ENG{implementation}\newline
  \textbullet~\ENG{Comprehensive} \ENG{conformance} \ENG{suite}\newline
  \textbullet~\ENG{Widely} \ENG{supported} \ENG{by} \ENG{LMS} \ENG{vendors}\newline
  \textbullet~\ENG{Enhanced} \ENG{data} \ENG{model} (4000-\ENG{char} \ENG{limits})
& \textbullet~\ENG{Still} \ENG{browser-dependent}\newline
  \textbullet~\ENG{No} \ENG{mobile} \ENG{optimization}\newline
  \textbullet~\ENG{No} \ENG{offline} \ENG{tracking}\newline
  \textbullet~\ENG{Limited} \ENG{to} \ENG{web} \ENG{context}\newline
  \textbullet~\textbf{\ENG{No} \ENG{updates} \ENG{since} 2009} \\

\end{longtable}
\normalsize

% \ENG{Historical} \ENG{context}
\begin{quote}
\footnotesize
\textbf{\ENG{Adoption} \ENG{Trends}:}
\begin{itemize}[noitemsep]
  \item \textbf{\ENG{SCORM} 1.2} (2001-2010): Κυρίαρχο \ENG{standard}, υποστηρίζεται από 95\%+ \ENG{LMS} \ENG{vendors}
  \item \textbf{\ENG{SCORM} 2004} (2004-\ENG{present}): Αργή υιοθέτηση λόγω πολυπλοκότητας, αλλά σήμερα \ENG{preferred} για νέα \ENG{projects}
  \item \textbf{\ENG{Legacy} \ENG{Impact}}: Πολλά \ENG{organizations} εξακολουθούν να χρησιμοποιούν \ENG{SCORM} 1.2 \ENG{content} λόγω \ENG{compatibility}
  \item \textbf{Μετά το 2010}: \ENG{xAPI} (2011) και \ENG{cmi5} (2016) εμφανίζονται ως \ENG{successors}, αλλά \ENG{SCORM} παραμένει ευρέως χρησιμοποιούμενο
\end{itemize}
\end{quote}



